\section{Related work}

\subsection{OOD detection}
The OOD detection community has explored a variety of techniques to underscore the distinctions between ID and OOD samples. These methods encompass classification-based~\cite{huang2021mos_confi,liang2018odin_confi,bendale2016towards_confi,devries2018learning_confi,hendrycks2016baseline_confi, sastry2020gram, tack2020csi, du2022vos}, density-based~\cite{zong2018deep_densi,abati2019latent_densi,nalisnick2018deep_densi,zisselman2020deep_densi,jiang2021revisiting_densi,pidhorskyi2018generative_densi,kirichenko2020normalizing_densi,sabokrou2018adversarially_densi}, and distance-based approaches~\cite{sun2022out_dist,lee2018maha_dist,ming2022cider_dist,chen2020boundary_dist,zaeemzadeh2021out_dist,van2020uncertainty_dist,techapanurak2020hyperparameter_dist,lu2023uncertainty_dist, sun2022knn}, with classification-based techniques generally outperforming the other types~\cite{yang2021generalized_survey}. In classification-based methods, the basic work of OOD detection starts with a simple and effective baseline: using the Maximum Softmax Probability (MSP)~\cite{hendrycks2016baseline_confi} to measure the probability that a certain sample is an ID sample.
On this basis, early approaches~\cite{liang2018odin_confi, hsu2020generalized, liu2020energy} focused on developing enhanced OOD indicators derived from neural network outputs. In addition, some researchers have proposed strategies involving OOD sample generation~\cite{lee2018training,du2022vos} and gradient-based~\cite{liang2018odin_confi} techniques. Among these, certain post-hoc methods~\cite{hendrycks2016msp, liang2018odin_confi,liu2020energy, sun2021react, sun2022dice, djurisic2022extremely, yu2023block,du2022vos} are notable for their simplicity and because they do not necessitate changes in the training process or objectives. This feature is particularly valuable for implementing OOD detection in real production environments, where the additional cost and complexity associated with retraining would be unacceptable.

The MSP method, initially presented by Hendrycks et al.~\cite{hendrycks2016msp}, was a formative step in post hoc OOD detection, using a neural network’s softmax output as a heuristic for distinguishing ID from OOD samples. Its straightforward application facilitated early adoption in OOD studies.
Despite MSP's influence, its limitations prompted further innovation, giving rise to the Energy method. This method, proposed by Liu et al.~\cite{liu2020energy}, refines the approach by assigning an energy score to network outputs, showing quantitative improvements over MSP with theoretical and empirical support.
Advancements in post hoc OOD detection have led to diverse methodological branches stemming from MSP~\cite{hendrycks2016baseline_confi} and Energy~\cite{liu2020energy} paradigms. LINe~\cite{ahn2023line} innovates by reducing neuron-induced noise through the calculation of Shapley values. Yu et al.~\cite{yu2023block} distinguish ID from OOD data by identifying neural network blocks with optimal differentiation based on the norms of their features. DICE~\cite{sun2022dice} improves discrimination by pruning weights in the fully connected layer according to the contribution units make during classification.
On the other end of the spectrum, entirely computation-free post hoc methods such as ReAct~\cite{sun2021react} and ASH~\cite{djurisic2022extremely} have shown promise. ReAct~\cite{sun2021react} investigates activations prior to the fully connected layer, applying rectification to suppress extreme activations that OOD data tend to trigger, thereby achieving refined detection outcomes. Similarly, ASH~\cite{djurisic2022extremely} prunes the activations inputted to the fully connected layer, but it achieves even more enhanced results compared to DICE~\cite{sun2022dice} by its selective pruning strategy. In this paper, our comparison mainly focuses on post hoc methods, since our method also belongs to this category.